\documentclass[12pt, a4paper]{article}

\usepackage[utf8]{inputenc}
\usepackage[T1]{fontenc}
\usepackage[brazil]{babel}
\usepackage{geometry}
\usepackage{amsmath}
\usepackage{amssymb}
\usepackage{booktabs}
\usepackage{hyperref}
\usepackage{xcolor}
\usepackage{listings}
\usepackage{graphicx}
\usepackage{float}
\usepackage{fancyhdr}
\usepackage{parskip}

\geometry{top=2.5cm, bottom=2.5cm, left=3cm, right=2.5cm, headheight=15pt}

\hypersetup{
    colorlinks=true,
    linkcolor=blue!70!black,
    urlcolor=blue!70!black,
    citecolor=blue!70!black
}

\definecolor{codegray}{rgb}{0.95,0.95,0.95}
\definecolor{codecomment}{rgb}{0.4,0.4,0.4}
\definecolor{codestring}{rgb}{0.13,0.54,0.13}
\definecolor{codekeyword}{rgb}{0.0,0.0,0.8}

\lstset{
    backgroundcolor=\color{codegray},
    commentstyle=\color{codecomment}\itshape,
    keywordstyle=\color{codekeyword}\bfseries,
    stringstyle=\color{codestring},
    basicstyle=\ttfamily\small,
    breakatwhitespace=false,
    breaklines=true,
    keepspaces=true,
    numbers=left,
    numbersep=8pt,
    numberstyle=\tiny\color{codecomment},
    showspaces=false,
    showstringspaces=false,
    showtabs=false,
    tabsize=4,
    frame=single,
    framerule=0.4pt,
    rulecolor=\color{gray!40}
}

\pagestyle{fancy}
\fancyhf{}
\fancyhead[L]{\small Home Credit Default Risk}
\fancyhead[R]{\small Valvitor Santos}
\fancyfoot[C]{\thepage}
\renewcommand{\headrulewidth}{0.4pt}

\title{
    \vspace{-1cm}
    {\Large \textbf{Home Credit Default Risk}} \\[0.5em]
    {\large Relat\'orio T\'ecnico --- Pipeline de Machine Learning\\
    para Predi\c{c}\~ao de Inadimpl\^encia}
}
\author{Valvitor Santos}
\date{\today}

% ──────────────────────────────────────────────────────────────
\begin{document}

\maketitle
\thispagestyle{fancy}

\begin{center}
\begin{minipage}{0.85\textwidth}
\small\textbf{Resumo.}
Este relat\'orio descreve o desenvolvimento de um pipeline completo de machine learning
para a competi\c{c}\~ao \textit{Home Credit Default Risk} do Kaggle. O problema consiste em
prever a probabilidade de inadimpl\^encia de clientes de cr\'edito com hist\'orico
financeiro limitado. A solu\c{c}\~ao emprega engenharia de features sobre sete tabelas
relacionais, modelo LightGBM com otimiza\c{c}\~ao de hiperpar\^ametros via Optuna e
valida\c{c}\~ao cruzada estratificada em 5 folds. A m\'etrica de avalia\c{c}\~ao \'e
a \textbf{AUC-ROC}.
\end{minipage}
\end{center}

\bigskip
\tableofcontents
\newpage

% ══════════════════════════════════════════════════════════════
\section{Introdu\c{c}\~ao e Contexto do Problema}

\subsection{Motiva\c{c}\~ao}

Institui\c{c}\~oes financeiras enfrentam o desafio de avaliar a capacidade de pagamento
de clientes com pouco ou nenhum hist\'orico de cr\'edito formal. A aus\^encia de dados
hist\'oricos em bureaus de cr\'edito convencionais leva \`a exclus\~ao financeira de
popula\c{c}\~oes que, de fato, s\~ao capazes de honrar seus compromissos. O Home Credit,
organiza\c{c}\~ao especializada em cr\'edito respons\'avel, busca usar fontes de dados
alternativas para tornar essa avalia\c{c}\~ao mais justa e precisa.

\subsection{Formula\c{c}\~ao do Problema}

O objetivo \'e construir um modelo de \textbf{classifica\c{c}\~ao bin\'aria} que, dado
o perfil de um solicitante, estime a probabilidade de inadimpl\^encia:

\begin{equation}
    \hat{p}_i = P(\text{TARGET}_i = 1 \mid \mathbf{x}_i)
\end{equation}

onde $\text{TARGET}_i = 1$ indica que o cliente $i$ n\~ao pagou o empr\'estimo, e
$\mathbf{x}_i$ \'e o vetor de features do cliente.

A m\'etrica de avalia\c{c}\~ao oficial \'e a \textbf{AUC-ROC} (\'Area sob a Curva ROC),
que mede a capacidade do modelo de separar inadimplentes de adimplentes,
independentemente do threshold de classifica\c{c}\~ao:

\begin{equation}
    \text{AUC} = \int_0^1 \text{TPR}\!\left(\text{FPR}^{-1}(t)\right) dt
\end{equation}

\subsection{Desafios Principais}

\begin{itemize}
    \item \textbf{Desbalanceamento de classes:} aproximadamente 8\% dos contratos resultam
          em inadimpl\^encia, criando forte vi\'es em favor da classe majorit\'aria.
    \item \textbf{Dados distribu\'idos em m\'ultiplas tabelas:} informa\c{c}\~oes relevantes
          est\~ao fragmentadas em 7 tabelas relacionais que precisam ser agregadas ao n\'ivel
          do solicitante.
    \item \textbf{Alta propor\c{c}\~ao de valores ausentes:} algumas features possuem at\'e
          60\% de nulos, exigindo tratamento cuidadoso.
    \item \textbf{Engenharia de features cr\'itica:} a performance depende fortemente das
          features derivadas, especialmente das tabelas auxiliares.
\end{itemize}

% ══════════════════════════════════════════════════════════════
\section{Conjunto de Dados}

\subsection{Estrutura Relacional}

Os dados s\~ao fornecidos em 7 arquivos CSV com a seguinte hierarquia:

\begin{table}[H]
\centering
\begin{tabular}{llrl}
\toprule
\textbf{Tabela} & \textbf{Chave} & \textbf{Linhas} & \textbf{Descri\c{c}\~ao} \\
\midrule
\texttt{application\_train/test} & \texttt{SK\_ID\_CURR} & 307.511 / 48.744 & Dados da solicita\c{c}\~ao atual \\
\texttt{bureau}                  & \texttt{SK\_ID\_CURR} & 1.716.428        & Cr\'editos em outros bancos \\
\texttt{bureau\_balance}         & \texttt{SK\_ID\_BUREAU}& 27.299.925      & Hist\'orico mensal do bureau \\
\texttt{previous\_application}   & \texttt{SK\_ID\_CURR} & 1.670.214        & Solicita\c{c}\~oes anteriores \\
\texttt{POS\_CASH\_balance}      & \texttt{SK\_ID\_CURR} & 10.001.358       & Saldos POS e empr\'estimos \\
\texttt{installments\_payments}  & \texttt{SK\_ID\_CURR} & 13.605.401       & Hist\'orico de parcelas \\
\texttt{credit\_card\_balance}   & \texttt{SK\_ID\_CURR} & 3.840.312        & Saldos de cart\~ao de cr\'edito \\
\bottomrule
\end{tabular}
\caption{Estrutura do conjunto de dados}
\end{table}

A tabela principal \'e a \^ancora do modelo; todas as demais s\~ao agregadas ao n\'ivel
\texttt{SK\_ID\_CURR} e unidas via \texttt{left join}.

\subsection{Vari\'avel Alvo}

\begin{itemize}
    \item $\texttt{TARGET} = 0$: cliente pagou normalmente (91,93\%)
    \item $\texttt{TARGET} = 1$: cliente inadimpliu (8,07\%)
\end{itemize}

% ══════════════════════════════════════════════════════════════
\section{An\'alise Explorat\'oria de Dados}

A EDA foi conduzida no notebook \texttt{01\_EDA.ipynb} com o objetivo de extrair insights
acion\'aveis --- n\~ao apenas estat\'isticas descritivas, mas decis\~oes de
pr\'e-processamento e engenharia de features diretamente informadas pelos dados.

\subsection{Anomalia em \texttt{DAYS\_EMPLOYED}}

A vari\'avel \texttt{DAYS\_EMPLOYED} representa h\'a quantos dias o cliente est\'a
empregado. Por\'em, 18.148 clientes apresentam o valor \textbf{365.243} --- um
placeholder para ``n\~ao empregado'', equivalente a $\approx$ 1000 anos de emprego.

\textbf{Decis\~ao:} substituir 365.243 por \texttt{NaN} e criar uma feature bin\'aria
\texttt{FLAG\_EMPLOYED\_ANOMALY} para preservar a informa\c{c}\~ao de que o dado estava
ausente.

\begin{lstlisting}[language=Python, caption={Tratamento da anomalia em DAYS\_EMPLOYED}]
df['FLAG_EMPLOYED_ANOMALY'] = (df['DAYS_EMPLOYED'] == 365243).astype(np.int8)
df['DAYS_EMPLOYED'] = df['DAYS_EMPLOYED'].replace(365243, np.nan)
\end{lstlisting}

\subsection{Features \texttt{EXT\_SOURCE}}

As vari\'aveis \texttt{EXT\_SOURCE\_1}, \texttt{EXT\_SOURCE\_2} e \texttt{EXT\_SOURCE\_3}
representam scores externos de cr\'edito provenientes de bureaus terceiros. A EDA revelou
que estas s\~ao as \textbf{features mais correlacionadas com a inadimpl\^encia} de todo
o dataset, tornando-as candidatas priorit\'arias para combina\c{c}\~oes n\~ao-lineares.

\texttt{EXT\_SOURCE\_1} possui 56\% de valores ausentes, exigindo a cria\c{c}\~ao de uma
flag adicional: \texttt{FLAG\_EXT\_SOURCE\_1\_MISSING}.

\subsection{Bloco Imobili\'ario}

Aproximadamente 20 vari\'aveis relacionadas ao im\'ovel do cliente (\texttt{APARTMENTS\_AVG},
\texttt{BASEMENTAREA\_AVG}, etc.) apresentam mais de 60\% de nulos e baixa correla\c{c}\~ao
com o target.

\textbf{Decis\~ao:} remover todas as colunas com mais de 60\% de valores ausentes,
exceto \texttt{EXT\_SOURCE}.

\subsection{Desbalanceamento de Classes}

Com 8,07\% de inadimpl\^encia, modelos n\~ao ajustados tendem a classificar quase todos
os clientes como adimplentes. \textbf{Decis\~ao:} usar \texttt{scale\_pos\_weight} no
LightGBM:

\begin{equation}
    \texttt{scale\_pos\_weight} = \frac{N_{\text{negativos}}}{N_{\text{positivos}}}
    = \frac{283{.}301}{24{.}210} \approx 11{,}70
\end{equation}

% ══════════════════════════════════════════════════════════════
\section{Engenharia de Features}

\subsection{Vis\~ao Geral}

A engenharia de features foi executada no notebook \texttt{02\_feature\_engineering.ipynb},
transformando as 7 tabelas brutas em um \'unico dataset processado com mais de 1.000
features. O processo segue quatro etapas:

\begin{enumerate}
    \item Features manuais na tabela principal
    \item Agrega\c{c}\~ao do bureau e bureau\_balance
    \item Agrega\c{c}\~ao das tabelas auxiliares (previous, POS, installments, credit card)
    \item Merge final, One-Hot Encoding e alinhamento
\end{enumerate}

\subsection{Features da Tabela Principal}

\subsubsection{Features Financeiras}

Raz\~oes financeiras que capturam a rela\c{c}\~ao entre cr\'edito, renda e parcelas:

\begin{table}[H]
\centering
\begin{tabular}{ll}
\toprule
\textbf{Feature} & \textbf{F\'ormula} \\
\midrule
\texttt{CREDIT\_INCOME\_RATIO}  & $\texttt{AMT\_CREDIT} / (\texttt{AMT\_INCOME\_TOTAL} + 1)$ \\
\texttt{ANNUITY\_INCOME\_RATIO} & $\texttt{AMT\_ANNUITY} / (\texttt{AMT\_INCOME\_TOTAL} + 1)$ \\
\texttt{ANNUITY\_CREDIT\_RATIO} & $\texttt{AMT\_ANNUITY} / (\texttt{AMT\_CREDIT} + 1)$ \\
\texttt{INCOME\_PER\_PERSON}    & $\texttt{AMT\_INCOME\_TOTAL} / (\texttt{CNT\_FAM\_MEMBERS} + 1)$ \\
\texttt{GOODS\_CREDIT\_RATIO}   & $\texttt{AMT\_GOODS\_PRICE} / (\texttt{AMT\_CREDIT} + 1)$ \\
\bottomrule
\end{tabular}
\caption{Features financeiras criadas}
\end{table}

O ac\'rescimo de 1 no denominador evita divis\~ao por zero sem distorcer valores
para denominadores grandes.

\subsubsection{Features Temporais}

\begin{itemize}
    \item \texttt{DAYS\_BIRTH\_YEARS}: idade em anos (mais interpret\'avel que dias negativos)
    \item \texttt{DAYS\_EMPLOYED\_YEARS}: tempo de emprego em anos
    \item \texttt{EMPLOYED\_TO\_BIRTH\_RATIO}: fra\c{c}\~ao da vida trabalhada
    \item \texttt{EMPLOYED\_TO\_AGE\_PCT}: propor\c{c}\~ao tempo de emprego / idade
\end{itemize}

\subsubsection{Combina\c{c}\~oes de \texttt{EXT\_SOURCE}}

Dado o alto poder preditivo dos scores externos, foram criadas m\'ultiplas combina\c{c}\~oes:

\begin{align}
    \texttt{EXT\_SOURCE\_MEAN} &= \tfrac{1}{3}(E_1 + E_2 + E_3) \\
    \texttt{EXT\_SOURCE\_STD}  &= \text{std}(E_1, E_2, E_3) \\
    \texttt{EXT\_SOURCE\_PROD} &= E_1 \cdot E_2 \cdot E_3 \\
    \texttt{EXT12} &= E_1 \cdot E_2, \quad
    \texttt{EXT13} = E_1 \cdot E_3, \quad
    \texttt{EXT23} = E_2 \cdot E_3
\end{align}

O produto captura intera\c{c}\~oes n\~ao-lineares: um score baixo em qualquer fonte
colapsa o produto, sinalizando risco concentrado.

\subsection{Agrega\c{c}\~ao das Tabelas Auxiliares}

Para cada tabela auxiliar, o processo \'e:

\begin{enumerate}
    \item Criar features manuais espec\'ificas para o dom\'inio da tabela
    \item Aplicar agrega\c{c}\~oes estat\'isticas para colunas num\'ericas:
          \texttt{count}, \texttt{mean}, \texttt{max}, \texttt{min}, \texttt{sum}
    \item Aplicar One-Hot Encoding + agrega\c{c}\~ao para colunas categ\'oricas:
          \texttt{sum}, \texttt{mean}
    \item Fazer \texttt{left join} ao n\'ivel \texttt{SK\_ID\_CURR}
\end{enumerate}

Este processo \'e implementado nas fun\c{c}\~oes reutiliz\'aveis
\texttt{agg\_numeric()} e \texttt{agg\_categorical()} em \texttt{src/features.py}.

\subsubsection{Bureau e Bureau Balance}

O \texttt{bureau\_balance} fornece o hist\'orico mensal de cada cr\'edito. As colunas
de status foram binarizadas:

\begin{lstlisting}[language=Python, caption={Engenharia de features no bureau\_balance}]
# C = pago, X = sem divida, 1-5 = meses de atraso crescente
bureau_balance['STATUS_GOOD'] = bureau_balance['STATUS'].isin(['C','X']).astype(np.int8)
bureau_balance['STATUS_LATE'] = bureau_balance['STATUS'].isin(['1','2','3','4','5']).astype(np.int8)
\end{lstlisting}

\subsubsection{Parcelas (\texttt{installments\_payments})}

\begin{itemize}
    \item \texttt{PAYMENT\_DIFF}: valor pago $-$ valor devido (positivo = pagou mais que o m\'inimo)
    \item \texttt{PAYMENT\_RATIO}: raz\~ao entre pago e devido
    \item \texttt{DAYS\_PAST\_DUE}: dias de atraso (zerado para pagamentos antecipados)
    \item \texttt{PAID\_ON\_TIME}: flag bin\'aria de pontualidade
\end{itemize}

\subsubsection{Cart\~ao de Cr\'edito (\texttt{credit\_card\_balance})}

\begin{equation}
    \texttt{UTILIZATION} = \frac{\texttt{AMT\_BALANCE}}{\texttt{AMT\_CREDIT\_LIMIT\_ACTUAL} + 1}
\end{equation}

Alta utiliza\c{c}\~ao do limite \'e um forte indicador de estresse financeiro.

\subsubsection{Aplica\c{c}\~oes Anteriores (\texttt{previous\_application})}

\begin{itemize}
    \item \texttt{APP\_CREDIT\_RATIO}: raz\~ao entre valor solicitado e aprovado ---
          captura se o cliente tipicamente pede mais do que recebe
    \item \texttt{DAYS\_DECISION\_YEARS}: tempo desde a \'ultima decis\~ao de cr\'edito
\end{itemize}

\subsection{Processamento Final}

Ap\'os o merge de todas as tabelas:

\begin{enumerate}
    \item Remo\c{c}\~ao de colunas com \textbf{mais de 80\% de nulos} p\'os-merge
    \item \textbf{One-Hot Encoding} para vari\'aveis categ\'oricas com menos de 10
          categorias \'unicas
    \item \textbf{Alinhamento} entre train e test: colunas ausentes no test s\~ao
          preenchidas com 0
    \item \textbf{Sanitiza\c{c}\~ao} de nomes de colunas: caracteres especiais
          substitu\'idos por \texttt{\_} (compatibilidade com LightGBM)
    \item Salvamento em \textbf{formato Parquet} (columnar, comprimido)
\end{enumerate}

O resultado final \'e um dataset com \textbf{1.042 features} e 307.511 linhas de treino.

% ══════════════════════════════════════════════════════════════
\section{Modelagem}

\subsection{Escolha do Algoritmo}

O \textbf{LightGBM} (Light Gradient Boosting Machine) foi escolhido pelos seguintes motivos:

\begin{itemize}
    \item \textbf{Efici\^encia com alta dimensionalidade:} os algoritmos GOSS e EFB
          permitem treinar rapidamente com mais de 1.000 features
    \item \textbf{Robustez a valores ausentes:} lida nativamente com NaN, aprendendo
          a melhor dire\c{c}\~ao de divis\~ao para dados ausentes
    \item \textbf{Performance em dados tabulares:} boosting de \'arvores \'e
          historicamente superior a redes neurais em dados tabulares com features
          de dom\'inio
    \item \textbf{Suporte a desbalanceamento:} par\^ametro \texttt{scale\_pos\_weight}
          ajusta o peso da classe minorit\'aria no c\'alculo do gradiente
\end{itemize}

\subsection{Otimiza\c{c}\~ao de Hiperpar\^ametros com Optuna}

O Optuna implementa o \textbf{TPE} (\textit{Tree-structured Parzen Estimator}), um
m\'etodo de otimiza\c{c}\~ao bayesiana que constr\'oi um modelo probabil\'istico do
espa\c{c}o de hiperpar\^ametros e sugere configura\c{c}\~oes com maior probabilidade
de melhoria.

O processo de tuning utilizou:
\begin{itemize}
    \item 20 \textit{trials}, cada um avaliado com 3-fold CV
    \item M\'etrica objetivo: AUC-ROC m\'edio dos 3 folds
    \item Early stopping de 30 rodadas por fold para acelerar a busca
\end{itemize}

\begin{table}[H]
\centering
\begin{tabular}{lll}
\toprule
\textbf{Hiperpar\^ametro} & \textbf{Espa\c{c}o de Busca} & \textbf{Tipo} \\
\midrule
\texttt{learning\_rate}       & $[0.02,\; 0.15]$          & log-uniform \\
\texttt{num\_leaves}          & $[20,\; 200]$             & inteiro \\
\texttt{max\_depth}           & $[4,\; 10]$               & inteiro \\
\texttt{min\_child\_samples}  & $[20,\; 300]$             & inteiro \\
\texttt{feature\_fraction}    & $[0.4,\; 0.9]$            & uniforme \\
\texttt{bagging\_fraction}    & $[0.4,\; 0.9]$            & uniforme \\
\texttt{bagging\_freq}        & $[1,\; 7]$                & inteiro \\
\texttt{reg\_alpha}           & $[10^{-4},\; 5.0]$        & log-uniform \\
\texttt{reg\_lambda}          & $[10^{-4},\; 5.0]$        & log-uniform \\
\bottomrule
\end{tabular}
\caption{Espa\c{c}o de busca do Optuna}
\end{table}

\subsection{Treinamento Final com 5-Fold CV}

Ap\'os identificar os melhores hiperpar\^ametros, o modelo final \'e treinado com
\textbf{StratifiedKFold} de 5 folds, que garante a preserva\c{c}\~ao da propor\c{c}\~ao
de inadimpl\^encia em cada fold:

\begin{lstlisting}[language=Python, caption={Treinamento com valida\c{c}\~ao cruzada estratificada}]
skf = StratifiedKFold(n_splits=5, shuffle=True, random_state=42)
oof_preds  = np.zeros(len(X))
test_preds = np.zeros(len(X_test))

for fold, (tr_idx, val_idx) in enumerate(skf.split(X, y)):
    model = lgb.LGBMClassifier(**BEST_PARAMS)
    model.fit(X.iloc[tr_idx], y.iloc[tr_idx],
              eval_set=[(X.iloc[val_idx], y.iloc[val_idx])],
              callbacks=[lgb.early_stopping(100)])
    oof_preds[val_idx]  = model.predict_proba(X.iloc[val_idx])[:, 1]
    test_preds         += model.predict_proba(X_test)[:, 1] / 5
\end{lstlisting}

As predi\c{c}\~oes de teste s\~ao a \textbf{m\'edia} das 5 predi\c{c}\~oes, reduzindo
vari\^ancia.

\subsection{M\'etrica OOF vs.\ M\'edia dos Folds}

O \textbf{AUC OOF} (Out-of-Fold) \'e calculado concatenando todas as predi\c{c}\~oes
out-of-fold e computando o AUC sobre o dataset completo:

\begin{equation}
    \text{AUC}_{\text{OOF}} = \text{AUC}\!\left(\mathbf{y}_{\text{train}},\;
    \hat{\mathbf{p}}_{\text{OOF}}\right)
\end{equation}

Isso \'e diferente da m\'edia das AUCs por fold:

\begin{equation}
    \overline{\text{AUC}}_{\text{folds}} = \frac{1}{K}
    \sum_{k=1}^{K} \text{AUC}\!\left(\mathbf{y}^{(k)}_{\text{val}},\;
    \hat{\mathbf{p}}^{(k)}_{\text{val}}\right)
\end{equation}

A AUC OOF \'e mais confi\'avel pois avalia o \textbf{ranking global} de todas as
predi\c{c}\~oes juntas. Folds individuais podem ter AUC mais alta por dados
estatisticamente mais sep\'araveis --- o OOF corrige esse vi\'es.

\subsection{Tracking com MLflow}

Todos os experimentos s\~ao registrados no MLflow:
\begin{itemize}
    \item Hiperpar\^ametros do melhor trial do Optuna
    \item AUC OOF final e AUC por fold
    \item N\'umero de features utilizadas
\end{itemize}

% ══════════════════════════════════════════════════════════════
\section{Boas Pr\'aticas e Reprodutibilidade}

\subsection{Estrutura de C\'odigo}

A separa\c{c}\~ao entre notebooks e c\'odigo reutiliz\'avel (\texttt{src/features.py})
evita duplica\c{c}\~ao. As fun\c{c}\~oes \texttt{agg\_numeric()} e
\texttt{agg\_categorical()} s\~ao gen\'ericas e funcionam para qualquer tabela,
recebendo apenas o DataFrame, a chave prim\'aria e o prefixo desejado.

\subsection{Gerenciamento de Mem\'oria}

Com datasets de dezenas de milh\~oes de linhas, o uso de mem\'oria \'e cr\'itico:

\begin{itemize}
    \item \textbf{Downcast de tipos:} inteiros e floats convertidos para o menor tipo
          que comporta seus valores (\texttt{int8}, \texttt{int16}, \texttt{float32})
    \item \textbf{Libera\c{c}\~ao expl\'icita:} ap\'os cada tabela auxiliar, o objeto
          \'e deletado e \texttt{gc.collect()} \'e chamado
    \item \textbf{Parquet:} armazenamento intermedi\'ario columnar e comprimido
\end{itemize}

\subsection{Preven\c{c}\~ao de Data Leakage}

O alinhamento de colunas entre train e test \'e feito via
\texttt{pd.DataFrame.align(join=`left')} para garantir que o test nunca possua features
derivadas de informa\c{c}\~oes do target. O Target Encoding, quando aplicado, usa
K-Fold para evitar vaza\c{c}\~ao.

% ══════════════════════════════════════════════════════════════
\section{Conclus\~ao e Pr\'oximos Passos}

\subsection{Conclus\~ao}

Este projeto implementa um pipeline end-to-end robusto para predi\c{c}\~ao de
inadimpl\^encia. Os pontos t\'ecnicos principais s\~ao:

\begin{enumerate}
    \item EDA orientada a decis\~oes, n\~ao apenas descritiva
    \item Engenharia de features fundamentada em dom\'inio financeiro
    \item Valida\c{c}\~ao cruzada estratificada para lidar com desbalanceamento
    \item Otimiza\c{c}\~ao bayesiana de hiperpar\^ametros com Optuna
    \item Tracking completo de experimentos com MLflow
\end{enumerate}

\subsection{Pr\'oximos Passos}

\begin{itemize}
    \item \textbf{Stacking/Blending:} combinar LightGBM com XGBoost e CatBoost
    \item \textbf{Target Encoding:} para vari\'aveis de alta cardinalidade como
          \texttt{ORGANIZATION\_TYPE}
    \item \textbf{SHAP values:} sele\c{c}\~ao de features baseada em importância de permuta\c{c}\~ao
    \item \textbf{Pseudo-labeling:} usar predi\c{c}\~oes de alta confiança no test
          para ampliar o treino
\end{itemize}

% ══════════════════════════════════════════════════════════════
\appendix
\section{Depend\^encias}

\begin{table}[H]
\centering
\begin{tabular}{ll}
\toprule
\textbf{Biblioteca} & \textbf{Vers\~ao} \\
\midrule
pandas      & 2.2.2  \\
numpy       & 1.26.4 \\
scikit-learn & 1.4.2 \\
lightgbm    & 4.x    \\
optuna      & 3.6.1  \\
mlflow      & 2.11.3 \\
pyarrow     & 15.0.2 \\
matplotlib  & 3.8.4  \\
seaborn     & 0.13.2 \\
\bottomrule
\end{tabular}
\caption{Depend\^encias principais do projeto}
\end{table}

\end{document}
